\setlength{\parindent}{1.25cm}

Primeiramente, expresso minha profunda gratidão a minha família, que sempre dedicou seus esforços à minha educação. Em especial, à minha mãe Mônica, que mesmo diante das circunstâncias mais adversas jamais mediu esforços para que seus filhos pudessem seguir seus caminhos e aproveitar cada oportunidade que a vida lhes oferecia. Ao meu pai George, meu eterno agradecimento por ter sido um exemplo vivo de superação, ensinando-me, pela própria trajetória, que sou capaz de vencer qualquer obstáculo que se apresente em meu caminho. 

Ao meu amor Geovanna, agradeço pelo amor e pelo apoio incondicional em todos os momentos desta jornada. Cada dia compartilhado, cada conversa e cada risada foram fundamentais para que eu chegasse até aqui. Sem a sua presença e dedicação, esta conquista não teria o mesmo significado. 

Aos meus amigos Carlos e César, agradeço por terem compartilhado comigo a trajetória de deixar o interior de Pernambuco em busca de novos horizontes em João Pessoa. As risadas, as conversas e os momentos vividos no apartamento 303 ficarão para sempre guardados em minha memória. 

A Isaías, um verdadeiro irmão que a Universidade Federal da Paraíba me presenteou, minha mais sincera gratidão. As inúmeras tardes e noites de diálogo, nas quais compartilhávamos e debatíamos ideias que transcendiam os limites acadêmicos. 

Agradeço a todos que, de alguma forma, cruzaram meu caminho e deixaram sua marca. Em especial, a Francisca, Lucélia e a toda a equipe ADM ProKids, que me ofereceram suporte precioso no momento mais difícil da minha vida. Estendo também meu profundo agradecimento a toda a equipe de oncologia e urologia do Hospital das Clínicas da Universidade Federal de Pernambuco, cuja dedicação e humanidade tornaram cada sessão de quimioterapia mais suportável. 

Por fim, ao meu orientador Yuri e ao meu coorientador Luiz, meu mais sincero e especial agradecimento. Ambos foram pilares fundamentais na construção deste trabalho. Estendo ainda minha gratidão a todo o corpo docente com quem tive o privilégio de aprender ao longo desta trajetória.